% chapters/testing/load_tests.tex
% Sección 6.4: Pruebas de carga (~0.5-1 página)
% ============================================================================
%
% GUÍA: Describir las pruebas de rendimiento bajo carga.
%
%   1. Herramienta utilizada:
%      TODO: Especificar (Locust, k6, Artillery, ab, wrk, etc.)
%
%   2. Escenarios de carga:
%      - Escenario 1: N usuarios concurrentes consultando exámenes.
%      - Escenario 2: Ingesta simultánea de M PDFs vía SFTP.
%      - Escenario 3: N correctores anotando exámenes en paralelo.
%      TODO: Definir los escenarios concretos.
%
%   3. Entorno de ejecución:
%      - Especificaciones del hardware de prueba.
%      - Configuración de Docker Compose (número de workers, etc.).
%
%   4. Resultados resumidos:
%      - Tiempos de respuesta (p50, p95, p99).
%      - Tasa de peticiones por segundo (RPS).
%      - Errores bajo carga.
%
%      \begin{table}{tab:load_results}{Resultados de pruebas de carga}
%        \begin{tabular}{l r r r r}
%          \textbf{Escenario} & \textbf{Usuarios} & \textbf{p50 (ms)} & \textbf{p95 (ms)} & \textbf{Errores (\%)} \\
%          \hline
%          Consulta exámenes   & -- & -- & -- & -- \\
%          Ingesta PDFs        & -- & -- & -- & -- \\
%          Corrección paralela & -- & -- & -- & -- \\
%        \end{tabular}
%      \end{table}
%
%   5. Análisis y conclusiones de rendimiento.
%      - Cuellos de botella identificados (si los hay).
%      - Optimizaciones aplicadas (si las hubo).
%
%   NOTA: Gráficas detalladas y logs completos van en \cref{ap:detailed_results}.
% ============================================================================

% TODO: Redactar la sección de pruebas de carga.
